% Iterated DS characters and the Delta5 Fourier boundary.

\providecommand{\Phi}{\Phi}
\providecommand{\Vir}{\mathrm{Vir}}
\providecommand{\sltwo}{\mathfrak{sl}_2}
\providecommand{\DS}{H^{0}_{DS}}
\providecommand{\Dfive}{D_5}
\providecommand{\kBKM}{\kappa_{\mathrm{BKM}}}
\providecommand{\wt}{\mathrm{wt}}
\providecommand{\schurch}{\operatorname{ch}}

\section{Iterated Drinfeld--Sokolov characters at the
\texorpdfstring{$\Delta_5$}{Delta5} Fourier boundary}
\label{sec:V24-Delta5-fourier-boundary}

\paragraph{Primitive denominator.}
Let
\[
  \mathcal{V}_{24}
  =
  \DS\!\bigl(L_{-2+1/22}(\sltwo)^{\otimes 22}\bigr)
\]
denote the iterated Drinfeld--Sokolov reduction attached to the
twenty-two-copy boundary construction.  The K3 Borcherds--Kac--Moody
branch is governed by the primitive Gritsenko--Nikulin denominator
$\Delta_5$.  The monic product normalisation used for scalar
Oberdieck--Pixton reduced-DT expressions is
\[
  \Dfive=64^{-1}\Delta_5 .
\]
The primitive Borcherds input is the Eichler--Zagier weak Jacobi form
$\phi_{0,1}^{\mathrm{EZ}}$ with
\[
  \phi_{0,1}^{\mathrm{EZ}}(\tau,z)\big|_{q^0}
  =
  y^{-1}+10+y .
\]
The K3 elliptic genus is
\[
  Z_{K3}=2\phi_{0,1}^{\mathrm{EZ}} .
\]
The primitive denominator uses $\phi_{0,1}^{\mathrm{EZ}}$ and therefore
\[
  c_1(0)=10,\qquad
  \kBKM(\Delta_5)=\wt(\Delta_5)=\frac{c_1(0)}2=5 .
\]

\paragraph{Square normalisation.}
The scalar square belongs to the partition-function convention:
\[
  \Phi_{10}^{\mathrm{un}}=\Delta_5^2,\qquad
  \Phi_{10}^{\mathrm{OP}}=\Dfive^2=4096^{-1}\Delta_5^2 .
\]
The quotient $\Phi_{10}/\Delta_5^2$ records square-normalisation data.
The primitive denominator, its weight, and the root multiplicities of
$\mathfrak g_{\Delta_5}$ remain governed by $\Delta_5$ and
$c_1(0)=10$.

\paragraph{Central charge normalisation.}
The Beem--Rastelli--Song four-dimensional to two-dimensional map gives
\[
  c_{2d}=-12c_{4d},\qquad c_{4d}=\frac{107}{6},
\]
and hence
\[
  c_{\mathcal V_{24}}=-214 .
\]
For the modularly shifted character
\[
  \chi^{\mathrm{mod}}_{\mathcal V_{24}}(\tau)
  =
  q^{-c_{\mathcal V_{24}}/24}
  \schurch_{\mathcal V_{24}}(\tau),
  \qquad
  \schurch_{\mathcal V_{24}}(\tau)
  =
  \operatorname{tr}_{\mathcal V_{24}}q^{L_0},
\]
one has
\[
  \chi^{\mathrm{mod}}_{\mathcal V_{24}}(\tau)
  =
  q^{107/12}\left(1+\sum_{n\geq 1}d_nq^n\right).
\]
Thus the genus-one graded character is obtained from
$\chi^{\mathrm{mod}}_{\mathcal V_{24}}$ by multiplying by
$q^{-107/12}$.

\paragraph{Humbert residue.}
Along the Humbert divisor, the regularised diagonal residue of the
squared primitive denominator is
\[
  \left[(2\pi iz)^2\Delta_5(\tau,z,\tau)^{-2}\right]_{z=0}
  =
  \eta(\tau)^{-48}
  =
  q^{-2}\prod_{m\geq 1}(1-q^m)^{-48}.
\]
Writing
\[
  E_{48}(q)=\prod_{m\geq 1}(1-q^m)^{-48}
  =\sum_{n\geq 0}e_nq^n,
\]
the coefficients follow from
\[
  (1-q^m)^{-48}
  =
  \sum_{j\geq 0}\binom{47+j}{47}q^{mj}.
\]
Through $q^{10}$ this gives
\[
\begin{array}{c|rrrrrrrrrrr}
n&0&1&2&3&4&5&6&7&8&9&10\\ \hline
e_n&
1&48&1224&21952&309876&3657312&37468928&
341773440&2826752418&21491641808&151810235136
\end{array}
\]
The same vector is obtained from the recurrence
\[
  n e_n
  =
  48\sum_{m=1}^{n}\sigma_1(m)e_{n-m},
  \qquad e_0=1.
\]

\paragraph{Character comparison.}
The coefficient comparison through $q^{10}$ is the genus-one statement
\[
  \schurch_{\mathcal V_{24}}(\tau)
  =
  E_{48}(q)
  \pmod{q^{11}},
\]
equivalently
\[
  q^{-107/12}\chi^{\mathrm{mod}}_{\mathcal V_{24}}(\tau)
  =
  q^2\left[(2\pi iz)^2\Delta_5(\tau,z,\tau)^{-2}\right]_{z=0}
  \pmod{q^{11}} .
\]
Its proof requires an exact BRST Euler--Poincare character formula for
the iterated Drinfeld--Sokolov reduction, the regularised Humbert
residue above, the Schur-sector identification of
$\mathcal V_{24}$, and equality with the vector
$(e_0,\ldots,e_{10})$.  A full genus-two identity requires a
superconformal index or a genus-two chiral block carrying the second
modular variable.

\paragraph{Borcherds constant.}
For the primitive K3 denominator,
\[
  \kBKM(\Delta_5)=5=\frac{10}{2}.
\]
For the squared scalar partition function,
\[
  \wt(\Delta_5^2)=10=2\kBKM(\Delta_5).
\]
The factor of two belongs to the physical square.  The primitive
Borcherds--Kac--Moody denominator and its root multiplicity convention
remain those of $\Delta_5$.

\paragraph{Source anchors.}
Gritsenko--Nikulin construct the primitive denominator $\Delta_5$ and
the monic product $\Dfive=64^{-1}\Delta_5$.  Borcherds'
singular-theta weight theorem gives $\wt(\Psi_f)=c_f(0)/2$.  The
$N=1$ K3 Jacobi input has $c_1(0)=10$ in the Eichler--Zagier
normalisation, while $Z_{K3}=2\phi_{0,1}^{\mathrm{EZ}}$ supplies the
doubled scalar square.  Beem--Rastelli and Song give the normalisation
$c_{2d}=-12c_{4d}$ and the value $c_{4d}=107/6$ for the Schur-sector
theory used here.

\paragraph{Drinfeld--Sokolov side.}
The denominator-side vector above is a direct product computation.  The
corresponding Drinfeld--Sokolov theorem is the independent character
computation through the same order, with the null-vector quotient and
the BRST grading convention fixed before coefficient comparison.
